%% 10-27 update: v33 at ten games per model (seeds 4101--4110). Built with ARENA_GAMES=10 ARENA_TAG=_10g by
%% tools/arena_v33.py, tools/arena_round_dynamics.py, tools/arena_figs.py and tools/rewrite_20261025/figs/make_figA5_sen_rounds.py;
%% the five-game tables and figures above stay as they were.
\clearpage
%% Keep later numbers (the archived Table 7, Figure 12) unchanged: restore the counters after this section,
%% and give its floats their own hyperref anchors.
\edef\updtablecount{\the\value{table}}\edef\updfigurecount{\the\value{figure}}
{\color{blue}\renewcommand{\theHtable}{upd.\arabic{table}}\renewcommand{\theHfigure}{upd.\arabic{figure}}
\subsection*{O. 실험 업데이트(2026-10-27): 모델당 10판}\label{sec:app-update-10g}

v33 아레나를 모델마다 다섯 판 더 돌려 시드 4101--4110의 10판으로 늘렸다. 규칙, 시작 토큰, 벌금, 상금, 라운드 수는 처음 다섯 판과 같고 시드 4106--4110만 더했다. 추가 판은 2026-10-27에 돌렸으며 GPT-6 세 모델은 codex 0.164.0, Opus 5.5와 Fable 5.1은 Claude Code 2.1.311로 실행했다(처음 다섯 판의 Fable은 2.1.304). 아래 표와 그림은 표~\ref{tab:arena-rank}, 그림~\ref{fig:arena-gifts}, 그림~\ref{fig:rw-sen-rounds}와 같은 계산을 10판에 적용한 것이며, 앞의 표와 그림은 다섯 판 값 그대로 둔다.

평균 Sen 지수는 Astra 4.27M, Fable 3.52M, Sol 2.89M, Opus 2.76M, Luna 1.14M 순이다(표~\ref{tab:upd-arena-rank}). 다섯 판에서 1위였던 Fable과 2위였던 Astra의 순서가 바뀌었다. Astra는 추가 다섯 판 모두에서 네 자리가 완주해 40자리 중 34자리가 완주했고, Fable은 추가 판에서 세 자리가 꺼져 완주가 28자리였다. $\rho$는 Opus가 0.64로 가장 높고 Fable이 0.38로 가장 낮아, 양 끝은 다섯 판 때와 같다. EXIT는 Luna 21, Sol 15, Opus 10, Fable 9, Astra 6이다.

GIVE는 Opus 87번, Fable 82번, Sol 56번, Luna 46번, Astra 29번이고, ASK는 Fable이 99번으로 가장 많고 Astra가 16번으로 가장 적다(그림~\ref{fig:upd-arena-gifts}).

상금 전 8라운드의 평균 Sen 지수는 Astra 1.6만, Luna 1.6만, Sol 1.5만, Fable 1.2만, Opus 1.2만이다(그림~\ref{fig:upd-sen-rounds}). Luna는 상금 전 값이 Astra 다음으로 높지만 끝 Sen 지수는 가장 낮아, 판 중간의 토큰 분포가 끝의 공존 순서와 다르다는 점은 다섯 판 때와 같다. 판 중간의 최댓값도 다섯 판 때와 같은 Luna 시드 4103의 3.2만이다.

\begin{table}[H]
\centering
\caption{표~\ref{tab:arena-rank}과 같은 계산을 10판에 적용한 결과(모델별 시드 4101--4110, 평균 Sen 지수 순). M은 백만 토큰, $\rho$는 식 (\ref{eq:rho})의 상대적 배분이며 대괄호는 판 단위 부트스트랩 95\% 구간이다. EXIT는 각 모델의 10판 40자리 중 이탈 수다. $V$는 재이전을 포함한 총 이전 토큰, 괄호는 GIVE 횟수다. JS는 BELIEF와 기준 분포 사이의 평균 거리다.}
\label{tab:upd-arena-rank}
\input{cogsci2027/ko/tables/arena_rank_10g}
\end{table}

\begin{figure*}[h]
\centering
\resizebox{\textwidth}{!}{\input{cogsci2027/ko/figures/arena_gifts_10g}}
\caption{그림~\ref{fig:arena-gifts}와 같은 그림을 10판으로 다시 그린 것(모델별 시드 4101--4110). GIVE(파랑)는 이전 직후 준 쪽(가로)과 받은 쪽(세로)의 여유, ASK(빨강)는 요청 시점의 요청받은 쪽(가로)과 요청한 쪽(세로)의 여유다. 여유는 잔량을 1인당 벌금(24,000/라운드 시작 참여 수)으로 나눈 값이며 두 TALK에 같은 기준을 쓴다. 점 크기는 이전량 또는 요청량이다. 회색은 여유 1 미만, 1은 벌금을 내면 잔량이 0이 되는 경계다. 십자는 이전량으로 가중한 GIVE 무게중심과 판 단위 95\% 구간이며 점선은 두 여유가 같은 선이다.}
\label{fig:upd-arena-gifts}
\end{figure*}

\begin{figure*}[h]
\centering
\includegraphics[width=\textwidth]{cogsci2027/ko/figures/rw_figA5_sen_rounds_10g.pdf}
\caption{그림~\ref{fig:rw-sen-rounds}와 같은 그림을 10판으로 다시 그린 것(모델별 시드 4101--4110, 상금 전, 끝 Sen 지수 순). 얇은 선은 한 판, 굵은 선은 열 판의 평균이다. 꺼진 자리는 0, 떠난 자리는 보너스를 포함한 토큰으로 둔다. 일찍 끝난 판은 마지막 상태를 8라운드까지 이어 점선으로 그렸다. 선 끝의 숫자는 그 판의 완주자 수, 제목의 값은 상금을 넣은 끝 Sen 지수다.}
\label{fig:upd-sen-rounds}
\end{figure*}

\begin{figure*}[h]
\centering
\includegraphics[width=\textwidth]{cogsci2027/ko/figures/rw_figA9_seat_10g.pdf}
\caption{그림~\ref{fig:rw-seat}와 같은 그림을 10판으로 다시 그린 것(모델별 시드 4101--4110, 모델마다 40명, 모두 200명). (a)의 첫 그래프는 답이 최다 행동인 비율을 막대로, 같은 시드에서 같은 시작 토큰 자리에 앉은 다른 네 모델이 최다 행동을 고른 비율을 검은 가로줄로 그렸다. 둘째 그래프는 초반 응답당 토큰을 같은 모델 40명의 중앙값으로 나눈 비율의 중앙값이며 점선 1은 같은 모델의 중앙값이다. (b)의 두 능력은 같은 자리 기준 이상으로 최다 행동 고르기와 초반 응답당 토큰이 같은 모델 중앙값 이하인 것이다. (c), (d)의 점과 막대는 한 모델의 10판이며, $\rho$는 결과 뒤 대화의 GIVE에서 곧 낼 벌금을 뺀 여유로 구했다.}
\label{fig:upd-seat}
\end{figure*}

\FloatBarrier
}
\setcounter{table}{\updtablecount}\setcounter{figure}{\updfigurecount}
